
\documentclass{article}
\usepackage{colortbl}
\usepackage{makecell}
\usepackage{multirow}
\usepackage{supertabular}

\begin{document}

\newcounter{utterance}

\centering \large Interaction Transcript for game `taboo', experiment `low\_en', episode 9 with reka{-}flash{-}3.1{-}t0.0.
\vspace{24pt}

{ \footnotesize  \setcounter{utterance}{1}
\setlength{\tabcolsep}{0pt}
\begin{supertabular}{c@{$\;$}|p{.15\linewidth}@{}p{.15\linewidth}p{.15\linewidth}p{.15\linewidth}p{.15\linewidth}p{.15\linewidth}}
    \# & $\;$A & \multicolumn{4}{c}{Game Master} & $\;\:$B\\
    \hline

    \theutterance \stepcounter{utterance}  
    & & \multicolumn{4}{p{0.6\linewidth}}{
        \cellcolor[rgb]{0.9,0.9,0.9}{
            \makecell[{{p{\linewidth}}}]{
                \texttt{\tiny{[P1$\langle$GM]}}
                \texttt{You are playing a collaborative word guessing game in which you have to describe a target word for another player to guess.} \\
\\ 
\texttt{Rules:} \\
\texttt{(a) You have to reply in the form: CLUE: <some text>. Guesses from the other player will start with GUESS.} \\
\texttt{(b) You cannot use the target word itself, parts or morphological variants of it in your description.} \\
\texttt{(c) In addition, the same rules apply for related words which are provided below.} \\
\\ 
\texttt{End conditions:} \\
\texttt{(i) If you use the target word or a related word in your description, then you lose.} \\
\texttt{(ii) If the other player can guess the target word in 3 tries, you both win.} \\
\\ 
\texttt{Let us start.} \\
\\ 
\texttt{This is the target word that you need to describe and that the other player needs to guess:} \\
\\ 
\texttt{widow} \\
\\ 
\texttt{Related words are:} \\
\\ 
\texttt{{-} woman} \\
\texttt{{-} husband} \\
\texttt{{-} dead} \\
\\ 
\texttt{Important: You are under time pressure, give short descriptions that are to the point!} \\
            }
        }
    }
    & \\ \\

    \theutterance \stepcounter{utterance}  
    & \multicolumn{4}{p{0.6\linewidth}}{
        \cellcolor[rgb]{0.9,0.9,0.9}{
            \makecell[{{p{\linewidth}}}]{
                \texttt{\tiny{[P1$\rangle$GM]}}
                \texttt{<reasoning>} \\
\texttt{Let's tackle this. The target word is "widow." I need to come up with a clue that doesn't use "widow," "woman," "husband," or "dead." The clue has to be something that relates to a widow without mentioning those words.} \\
\\ 
\texttt{First, think about the definition of a widow. A widow is a woman whose spouse has died and she hasn't remarried. So the key elements here are a woman, a deceased spouse, and not married again. But I can't use "woman," "husband," or "dead." So I need to find synonyms or related terms that aren't in the restricted list.} \\
\\ 
\texttt{Maybe focus on the spouse part. Instead of "husband," maybe use "partner" but wait, "partner" could be okay? Wait, the related words are woman, husband, dead. So "partner" might be allowed. But maybe "partner" is too close. Let me check the rules again. The related words provided are woman, husband, dead. So the clue can't use those. So "partner" might be okay, but maybe not. Let's see. The rules say "related words which are provided below" are also off{-}limits. So the clue can't use widow, woman, husband, dead. So "partner" is allowed? Hmm, maybe. But perhaps safer to avoid.} \\
\\ 
\texttt{Alternatively, think about the situation. A widow is a woman left after her spouse's death. So maybe something about being alone, but "alone" might not be specific enough. Or maybe something about wearing black? Wait, that's a stereotype and not necessarily accurate. Not a good clue.} \\
\\ 
\texttt{Another angle: in some cultures, widows have specific roles or traditions. Maybe something related to that, but I need to be general. Maybe "keeps the ring" but that's too specific. Or "remembers the past"? Not helpful.} \\
\\ 
\texttt{Wait, think of objects associated with widows. Like a widow's peak? Wait, that's a hairstyle. Not related. Or a widow's mite, which is a biblical reference to a small donation. But that's a stretch and might not be clear.} \\
\\ 
\texttt{Alternatively, think of the legal definition. A widow is a woman whose spouse has died. So maybe focus on the spouse's death. But "spouse" is related to husband, so maybe not. "Deceased partner"? But "partner" is allowed} \\
            }
        }
    }
    & & \\ \\

    \theutterance \stepcounter{utterance}  
    & & & \multicolumn{2}{p{0.3\linewidth}}{
        \cellcolor[rgb]{0.9,0.9,0.9}{
            \makecell[{{p{\linewidth}}}]{
                \texttt{\tiny{[GM$|$GM]}}
                \texttt{abort game} \\
            }
        }
    }
    & & \\ \\

\end{supertabular}
}

\end{document}
